\textbf{结论名称：}存在性\(\Rightarrow\)最大值/最小值转化

\textbf{适用条件：}函数 \( f(x) \) 在区间 \( D \) 上有最大值和最小值（连续函数必满足）；讨论存在 \( x\in D \) 使得不等式成立的问题。

\textbf{变量说明：} \( D \) 为给定区间，\( f(x) \) 为定义在 \( D \) 上的连续函数，\( a \) 为常数参数。

\textbf{核心公式：}
\[
\boxed{\exists\, x\in D,\ f(x)\ge a \iff \max_{x\in D}f(x)\ge a}
\]
\[
\boxed{\exists\, x\in D,\ f(x)\le a \iff \min_{x\in D}f(x)\le a}
\]

\textbf{使用场景：} 含参数的方程或不等式存在解时，求参数取值范围。将存在量词直接消去，转化为最值与常数的大小比较，避免对原不等式进行复杂讨论。

\textbf{简单例子：} 设 \( f(x)=x^2-2x,\ x\in[0,2] \)。若存在 \( x\in[0,2] \) 使 \( f(x)\ge m \) 成立，求 \( m \) 的取值范围。
\[
\max_{x\in[0,2]}f(x)=f(0)=f(2)=0,\quad \text{故 } m\le 0.
\]
验证：取 \( m=-1 \)，存在 \( x=1 \) 使 \( f(1)=-1\ge -1 \) 成立；取 \( m=1 \)，对任意 \( x \) 均有 \( f(x)\le 0<1 \)，不存在，符合结论。